% Драфт секции 3.3.1: Постановка задачи максимизации активации нейрона
% Источник: parts/neurotrain/main.tex, строки 101–167
% Статус: черновик для ревью

\link{Связка: методы, описанные в предыдущих секциях, анализируют корреляции между зарегистрированной активностью нейронов и наблюдаемым поведением или свойствами среды. Однако задачу интерпретации функции нейрона можно поставить принципиально иначе: не искать корреляции с уже зарегистрированными стимулами, а \textit{создать} оптимальный стимул, вызывающий максимальный отклик данного нейрона. Такой подход --- максимизация активации --- позволяет перейти от пассивного наблюдения к целенаправленному зондированию нейронных репрезентаций.}

\subsection{Постановка задачи максимизации активации нейрона}\label{sec:mango_problem}

Нейроны мозга селективно реагируют на определённые свойства стимулов: нейроны первичной зрительной коры специализируются на простых признаках видимых изображений, таких как направление движения~\cite{HubelWiesel} или цвет~\cite{BiPOLES}, тогда как нейроны высших зрительных областей отвечают на более сложные стимулы, например на наличие лица в поле зрения~\cite{QuianQuiroga2023}. Выявление такой специализации включает итеративное изменение свойств предъявляемого стимула с целью поиска тех, которые вызывают наиболее интенсивный отклик~\cite{ponce2019evolving, bardon2022face}. Данный класс подходов известен как \textbf{максимизация активации} (МА).

С математической точки зрения активация нейрона рассматривается как функция входного стимула. Наиболее активирующее изображение (MEI, \textit{most exciting image}~\cite{Inception-what-excites-neurons-most}) определяется как
\begin{equation}\label{eq:mei_def}
    \textrm{MEI}_i \overset{\textrm{def}}{=} \underset{x \in S}{\arg\max}\; A_i(x),
\end{equation}
где $S$~--- пространство стимулов, $A_i(x)$~--- функция активации $i$-го нейрона. На практике $A_i(x)$ всегда измеряется с шумом, а количество доступных итераций ограничено условиями эксперимента, поэтому необходимо определить MEI как можно более эффективно --- за минимальное число предъявлений стимула~\cite{ponce2019evolving}.

Генеративные модели, такие как GAN~\cite{GANpaper}, позволяют параметризовать пространство стимулов $S$ компактным латентным вектором, существенно сокращая размерность оптимизационной задачи. Эффективность этого подхода была показана в работах XDream~\cite{XDream} и Inception Loops~\cite{Inception-what-excites-neurons-most}, в том числе при идентификации MEI нейронов зрительной коры макак \textit{in vivo}~\cite{ponce2019evolving, bardon2022face}.

При изучении биологической нейронной сети внутренняя структура системы недоступна, и активация нейрона представляет собой функцию типа <<чёрного ящика>> --- без явного доступа к градиенту. Аппроксимация градиента конечными разностями в таких условиях ненадёжна из-за шума измерений и потенциально сложного ландшафта функции активации. Это делает \textbf{безградиентные методы} оптимизации естественным выбором для решения задачи МА в реальных системах. В работе~\cite{TTOpt-paper} было показано, что подход на основе разложения тензорного поезда (TT-разложения~\cite{TT-Oseledets}) обладает значительными преимуществами перед альтернативными безградиентными методами (в том числе генетическими и эволюционными алгоритмами) по точности и скорости сходимости при фиксированном числе итераций. TT-разложение позволяет представить тензор в компактной малоранговой форме, линейной по размерности, и эффективно решать задачу оптимизации без вычисления производных~\cite{TTOpt-paper, PROTES-Chertkov-paper, TT-optimization-One-More-Paper-from-Andrei}.

Импульсные нейронные сети, описанные в разделе~\ref{sec: lit_ann_repr}, представляют собой более биологически реалистичную модель нейронных сетей мозга по сравнению с классическими сетями прямого распространения~\cite{SNN-biological-relevance, SNNs-what-are-overview}. В отличие от последних, активность элементов импульсных сетей разворачивается во времени --- было показано, что временная структура импульсной последовательности играет ключевую роль в функционировании мозга~\cite{Panzeri2001, AndradeTalavera2023}. Внутренние представления данных в импульсных нейронных сетях изучались в~\cite{SNN-neural-representations-2023}, однако задача поиска MEI для нейронов таких сетей безградиентными методами ранее не ставилась.

В настоящей работе предложен подход к безградиентной максимизации активации нейронов импульсных нейронных сетей на основе TT-оптимизации в латентном пространстве генеративной модели. Разработан программный фреймворк MANGO (\textit{Maximization of Activation of Neurons through Gradient-free Optimization}), объединяющий генератор стимулов, оптимизатор на основе TT-разложения и средства анализа полученных MEI. Насколько нам известно, применение безградиентных методов оптимизации стимулов к импульсным нейронным сетям реализовано впервые~\cite{Pospelov2025}.
